% Options for packages loaded elsewhere
\PassOptionsToPackage{unicode}{hyperref}
\PassOptionsToPackage{hyphens}{url}
%
\documentclass[
  ignorenonframetext,
]{beamer}
\usepackage{pgfpages}
\setbeamertemplate{caption}[numbered]
\setbeamertemplate{caption label separator}{: }
\setbeamercolor{caption name}{fg=normal text.fg}
\beamertemplatenavigationsymbolsempty
% Prevent slide breaks in the middle of a paragraph
\widowpenalties 1 10000
\raggedbottom
\setbeamertemplate{part page}{
  \centering
  \begin{beamercolorbox}[sep=16pt,center]{part title}
    \usebeamerfont{part title}\insertpart\par
  \end{beamercolorbox}
}
\setbeamertemplate{section page}{
  \centering
  \begin{beamercolorbox}[sep=12pt,center]{section title}
    \usebeamerfont{section title}\insertsection\par
  \end{beamercolorbox}
}
\setbeamertemplate{subsection page}{
  \centering
  \begin{beamercolorbox}[sep=8pt,center]{subsection title}
    \usebeamerfont{subsection title}\insertsubsection\par
  \end{beamercolorbox}
}
\AtBeginPart{
  \frame{\partpage}
}
\AtBeginSection{
  \ifbibliography
  \else
    \frame{\sectionpage}
  \fi
}
\AtBeginSubsection{
  \frame{\subsectionpage}
}

\usepackage{amsmath,amssymb}
\usepackage{iftex}
\ifPDFTeX
  \usepackage[T1]{fontenc}
  \usepackage[utf8]{inputenc}
  \usepackage{textcomp} % provide euro and other symbols
\else % if luatex or xetex
  \usepackage{unicode-math}
  \defaultfontfeatures{Scale=MatchLowercase}
  \defaultfontfeatures[\rmfamily]{Ligatures=TeX,Scale=1}
\fi
\usepackage{lmodern}
\usetheme[]{madrid}
\usecolortheme{default}
\usefonttheme{default}
\ifPDFTeX\else  
    % xetex/luatex font selection
\fi
% Use upquote if available, for straight quotes in verbatim environments
\IfFileExists{upquote.sty}{\usepackage{upquote}}{}
\IfFileExists{microtype.sty}{% use microtype if available
  \usepackage[]{microtype}
  \UseMicrotypeSet[protrusion]{basicmath} % disable protrusion for tt fonts
}{}
\makeatletter
\@ifundefined{KOMAClassName}{% if non-KOMA class
  \IfFileExists{parskip.sty}{%
    \usepackage{parskip}
  }{% else
    \setlength{\parindent}{0pt}
    \setlength{\parskip}{6pt plus 2pt minus 1pt}}
}{% if KOMA class
  \KOMAoptions{parskip=half}}
\makeatother
\usepackage{xcolor}
\newif\ifbibliography
\setlength{\emergencystretch}{3em} % prevent overfull lines
\setcounter{secnumdepth}{-\maxdimen} % remove section numbering


\providecommand{\tightlist}{%
  \setlength{\itemsep}{0pt}\setlength{\parskip}{0pt}}\usepackage{longtable,booktabs,array}
\usepackage{calc} % for calculating minipage widths
\usepackage{caption}
% Make caption package work with longtable
\makeatletter
\def\fnum@table{\tablename~\thetable}
\makeatother
\usepackage{graphicx}
\makeatletter
\newsavebox\pandoc@box
\newcommand*\pandocbounded[1]{% scales image to fit in text height/width
  \sbox\pandoc@box{#1}%
  \Gscale@div\@tempa{\textheight}{\dimexpr\ht\pandoc@box+\dp\pandoc@box\relax}%
  \Gscale@div\@tempb{\linewidth}{\wd\pandoc@box}%
  \ifdim\@tempb\p@<\@tempa\p@\let\@tempa\@tempb\fi% select the smaller of both
  \ifdim\@tempa\p@<\p@\scalebox{\@tempa}{\usebox\pandoc@box}%
  \else\usebox{\pandoc@box}%
  \fi%
}
% Set default figure placement to htbp
\def\fps@figure{htbp}
\makeatother

\AtBeginSubsection{}
\usepackage{amsfonts}
\setbeamertemplate{section page}{ \begin{centering} \begin{beamercolorbox}[sep=12pt,center,shadow=true,rounded=true]{section title} \usebeamerfont{section title}\insertsection\par \end{beamercolorbox} \end{centering}}
\setbeamertemplate{subsection page}{}
\setbeamertemplate{headline}{ \begin{beamercolorbox}[ht=2.5ex,dp=1.125ex,leftskip=.3cm,rightskip=.3cm]{section in head/foot} \usebeamerfont{section in head/foot}\insertsectionhead \hfill \usebeamerfont{subsection in head/foot}\insertsubsectionhead \end{beamercolorbox}}
\makeatletter
\@ifpackageloaded{caption}{}{\usepackage{caption}}
\AtBeginDocument{%
\ifdefined\contentsname
  \renewcommand*\contentsname{Table of contents}
\else
  \newcommand\contentsname{Table of contents}
\fi
\ifdefined\listfigurename
  \renewcommand*\listfigurename{List of Figures}
\else
  \newcommand\listfigurename{List of Figures}
\fi
\ifdefined\listtablename
  \renewcommand*\listtablename{List of Tables}
\else
  \newcommand\listtablename{List of Tables}
\fi
\ifdefined\figurename
  \renewcommand*\figurename{Figure}
\else
  \newcommand\figurename{Figure}
\fi
\ifdefined\tablename
  \renewcommand*\tablename{Table}
\else
  \newcommand\tablename{Table}
\fi
}
\@ifpackageloaded{float}{}{\usepackage{float}}
\floatstyle{ruled}
\@ifundefined{c@chapter}{\newfloat{codelisting}{h}{lop}}{\newfloat{codelisting}{h}{lop}[chapter]}
\floatname{codelisting}{Listing}
\newcommand*\listoflistings{\listof{codelisting}{List of Listings}}
\makeatother
\makeatletter
\makeatother
\makeatletter
\@ifpackageloaded{caption}{}{\usepackage{caption}}
\@ifpackageloaded{subcaption}{}{\usepackage{subcaption}}
\makeatother

\usepackage{bookmark}

\IfFileExists{xurl.sty}{\usepackage{xurl}}{} % add URL line breaks if available
\urlstyle{same} % disable monospaced font for URLs
\hypersetup{
  pdftitle={Markov Random Fields with Geographic Seeds},
  pdfauthor={Wesley Cabral Alves},
  hidelinks,
  pdfcreator={LaTeX via pandoc}}


\title{Markov Random Fields with Geographic Seeds}
\subtitle{Bayesian Estimation on Hierarchical Hidden Field Model}
\author{Wesley Cabral Alves}
\date{2026-08-30}
\institute{University of Campinas}

\begin{document}
\frame{\titlepage}


\section{1. Introduction}\label{introduction}

\subsection{Spatial Motivation}\label{spatial-motivation}

\begin{frame}{The Potts Model}
\phantomsection\label{the-potts-model}
\begin{itemize}
\tightlist
\item
  Spatial models on regular lattices (pixels) often employ
  \textbf{Markov Random Fields (MRF)}.
\item
  The \textbf{Potts Model} is the classical standard for capturing local
  neighborhood dependence.
\end{itemize}
\end{frame}

\begin{frame}{The Geographic Limitation}
\phantomsection\label{the-geographic-limitation}
\begin{itemize}
\tightlist
\item
  Many real images and textures possess geographic patterns we call
  \textbf{``seeds''} that may attract or repel colors.
\item
  \textbf{Our Proposal:} Incorporate geographic seeds that exert radial
  influence on the classes.
\item
  \textbf{Objective:} Propose a reparameterized model that describe the
  inclusion of centroids/seeds, and validate its inference under a
  hidden model.
\end{itemize}
\end{frame}

\section{2. The Proposal}\label{the-proposal}

\subsection{Mathematical Formulation}\label{mathematical-formulation}

\begin{frame}{Potts as is}
\phantomsection\label{potts-as-is}
\begin{itemize}
\tightlist
\item
  \textbf{Energy Function:} For the traditional Potts Model, we define a
  spatial field \(Z\) where the local probability of pixel \(i\) taking
  class \(k\) is proportional to \(exp(\eta_i(k))\):
  \[\eta_i(k) = \alpha \sum_{j \in N(i)} \mathbb{1}(Z_j = k) + \beta \mathbb{1}(k = 0)\]
\item
  \textbf{Parameters Meaning:}

  \begin{itemize}
  \tightlist
  \item
    \textbf{\(N(i)\):} Pixel \(i\) neighbors.
  \item
    \textbf{\(\alpha\) (Uniformity):} Local attraction between
    neighboring pixels (smoothing).
  \item
    \textbf{\(\beta\) (Global Abundance):} Force favoring/penalizing the
    background (class \(0\)).
  \end{itemize}
\end{itemize}
\end{frame}

\begin{frame}{Potts Reparameterized with Seeds}
\phantomsection\label{potts-reparameterized-with-seeds}
\begin{itemize}
\tightlist
\item
  \textbf{Enhanced Local Energy Function:} To address the isotropic
  limitation of the Potts model, we introduce a few elements to handle
  the geographic centrality bias in our local energy funtion:
  \[\eta_i(k) = \alpha \sum_{j \in N(i)} 1(Z_j = k) + \beta 1(k = 0) + \sum_s \frac{\delta_s}{1 + d(i, s)} 1(c_s = k)\]
\item
  \textbf{Parameters Meaning:}

  \begin{itemize}
  \tightlist
  \item
    \textbf{\(\alpha\) (Uniformity):} Local attraction between
    neighboring pixels (smoothing).
  \item
    \textbf{\(\beta\) (Global Abundance):} Force favoring/penalizing the
    background (class \(0\)).
  \item
    \textbf{\(\delta_s\) (Centrality):} Radial attraction/repulsion of
    the seed \(s\).
  \item
    \textbf{\(d(i, s)\):} Distance between pixel \(i\) and seed \(s\).
  \item
    \textbf{\(c_s\):} Class/color related to the seed.
  \end{itemize}
\end{itemize}
\end{frame}

\begin{frame}{Potts Reparameterized with Seeds}
\phantomsection\label{potts-reparameterized-with-seeds-1}
\begin{itemize}
\item
  \textbf{Enhanced Local Energy Function:}
  \[\eta_i(k) = \alpha \sum_{j \in N(i)} 1(Z_j = k) + \beta 1(k = 0) + \sum_s \frac{\delta_s}{1 + d(i, s)} 1(c_s = k)\]
\item
  \textbf{Enhanced Global Energy Function:}
  \[U(\omega) = \alpha  \sum_{\{i,j\} \in E} \sum_{k \in S} s_{i,k}s_{j,k} + \beta 1(k = 0) + \sum_s \frac{\delta_s}{1 + d(i, s)} 1(c_s = k)\]
\end{itemize}
\end{frame}

\begin{frame}{The Threshold of the ``Area of Influence''}
\phantomsection\label{the-threshold-of-the-area-of-influence}
\begin{itemize}
\tightlist
\item
  \textbf{Equilibrium of Forces:} The local seed attraction competes
  against the global abundance of the background class.
\item
  \textbf{Null Boundary:} Under our model, ignoring neighborhood
  smoothing, they cancel out when:
  \[\eta_i(0)= \eta_i(k) \implies d_{i,s} = \frac{\delta_s}{\beta}-1\]
\item
  \textbf{Geographic Transition:}

  \begin{itemize}
  \tightlist
  \item
    \textbf{Within radius \(d_i\):} The seed attraction dominates the
    local probability.
  \item
    \textbf{Beyond radius \(d_i\):} The class 0 dominates.
  \item
    \textbf{Smoothing effect:} Depending on \(\alpha\) value, pixels can
    be affected more by its naighbors than by the radius rule.
  \end{itemize}
\end{itemize}
\end{frame}

\begin{frame}{The Theoretical Bottleneck: Intractable Likelihood}
\phantomsection\label{the-theoretical-bottleneck-intractable-likelihood}
\begin{itemize}
\tightlist
\item
  \textbf{The Gibbs Measure:} The global probability of an image
  configuration \(\omega\) is governed by:
  \[P(\omega) = \frac{1}{Z(\theta)} \exp\{-U(\omega)\}\]
\item
  \textbf{The Computational Bottleneck:} The term \(Z(\theta)\) is the
  \textbf{normalizing constant} (partition function) summed over all
  lattice configurations \(\Omega\):
  \[Z(\theta) = \sum_{\omega \in \Omega} \exp\{-U(\omega)\}\]
\item
  For a small \(50 \times 50\) grid and \(K = 3\) colors, this requires
  \(3^{2500}\) terms, rendering standard Maximum Likelihood impossible.
\end{itemize}
\end{frame}

\section{3. Hidden \& Hierarchical
Models}\label{hidden-hierarchical-models}

\subsection{Hidden Markov Fields}\label{hidden-markov-fields}

\begin{frame}{The Major Evolution: The Hidden Model (HMRF)}
\phantomsection\label{the-major-evolution-the-hidden-model-hmrf}
\begin{itemize}
\tightlist
\item
  \textbf{The Latent Field:} In segmentation practice, the label field
  \(Z\) (the true texture) is \textbf{latent/hidden}.
\item
  \textbf{Observed Intensities:} We only observe the noisy, continuous
  intensities \(Y\).
\item
  \textbf{Observation Equation (Emission):} Conditionally on the latent
  label \(Z_i = k\), the intensity \(Y_i\) is Normal:
  \[Y_i \mid Z_i = k \sim N(\mu_k, \sigma^2_k)\]
\item
  \textbf{Normal Mixture:} Marginalizing over \(Z_i\), \(Y_i\) forms a
  spatial mixture of Gaussians:
  \[f(y_i) = \sum_{k=0}^{K-1} P(Z_i = k) \phi(y_i; \mu_k, \sigma^2_k)\]
\end{itemize}
\end{frame}

\subsection{Hierarchical Framework}\label{hierarchical-framework}

\begin{frame}{Hierarchical Bayesian Formulation}
\phantomsection\label{hierarchical-bayesian-formulation}
\begin{itemize}
\item
  \textbf{Full Joint Distribution:} Treating \(Z\), parameters
  \(\theta\), and seed coordinates \(S\) as random variables: \[
  \begin{aligned}
  P(Y, Z, \theta) = P(Y \mid Z, \mu, \sigma^2) P(Z \mid \alpha, \beta, \delta, S) \times  \\
  P(\alpha)P(\beta)\prod_{j=1}^M P(s_j)P(\delta_j)
  \prod_{k=0}^{K-1} P(\mu_k)P(\sigma^2_k)
  \end{aligned}
  \]
\item
  \textbf{Prior Distributions:}

  \begin{itemize}
  \tightlist
  \item
    \textbf{Field (Gamma):}
    \(\alpha \sim \text{Gamma}(a_\alpha, b_\alpha)\),
    \(\beta \sim \text{Gamma}(a_\beta, b_\beta)\).
  \item
    \textbf{Means (Normal):} \(\mu_k \sim N(m_k, \tau^2_k)\) (anchored
    via K-means).
  \item
    \textbf{Variances (Inverse-Gamma):}
    \(\sigma^2_k \sim IG(a_\sigma, b_\sigma)\).
  \item
    \textbf{Geography:} Seed position \(s_j \sim \text{Unif}(L)\)
    discretized on the grid.
  \end{itemize}
\end{itemize}
\end{frame}

\section{3. Estimation \& Algorithms}\label{estimation-algorithms}

\subsection{MCMC Estimation}\label{mcmc-estimation}

\begin{frame}{The Metropolis-within-Gibbs Algorithm}
\phantomsection\label{the-metropolis-within-gibbs-algorithm}
\begin{enumerate}
\tightlist
\item
  \textbf{Step C (Gibbs - Latent Field):} Sample \(Z_i^{(t)}\)
  conditioned on its neighbors and the observed intensity \(Y_i\).
\item
  \textbf{Step D (Gibbs - Noise):} Sample new \((\mu_k, \sigma^2_k)\)
  directly from closed-form analytical distributions.
\item
  \textbf{Step A (Metropolis-Hastings - Field):} Propose
  \((\alpha, \beta, \delta)\) via a random walk, evaluated over the
  \textbf{Pseudolikelihood (PL)} or user-chosen prioris.
\item
  \textbf{Step B (Metropolis-Hastings - Seeds):} Propose discrete
  1-pixel steps for the coordinates \((x,y)\) of the seeds.
\end{enumerate}

\begin{itemize}
\tightlist
\item
  \textbf{The Local Gibbs Trick:} Global normalizing constant
  \(Z(\theta)\) cancels out locally.
  \[P(Z_i = k \mid Y_i, Z_{-i}, \theta) = \frac{\phi(y_i; \mu_k, \sigma^2_k) \exp\{\eta_i(k)\}}{\sum_{\ell} \phi(y_i; \mu_\ell, \sigma^2_\ell) \exp\{\eta_i(\ell)\}}\]
\end{itemize}
\end{frame}

\begin{frame}{Full Conditionals and Conjugacy}
\phantomsection\label{full-conditionals-and-conjugacy}
\begin{itemize}
\tightlist
\item
  \textbf{Updated Posterior Mean (\(\mu_k\)):} The full conditional
  under conjugate priors is a Normal:
  \[\mu_{\text{post}} = \frac{\frac{m_k}{\tau^2_k} + \frac{\sum_{i: Z_i=k} y_i}{\sigma^2_k}}{\frac{1}{\tau^2_k} + \frac{n_k}{\sigma^2_k}}\]
\item
  \textbf{Updated Posterior Variance (\(\sigma^2_k\)):} Follows an
  Inverse-Gamma: \[a_{\text{post}} = a_\sigma + \frac{n_k}{2}\]
  \[b_{\text{post}} = b_\sigma + \frac{1}{2} \sum_{i: Z_i=k} (y_i - \mu_k)^2\]
\end{itemize}
\end{frame}

\subsection{Computational Efficiency}\label{computational-efficiency}

\begin{frame}{Code Engineering \& Stability Challenges}
\phantomsection\label{code-engineering-stability-challenges}
\begin{itemize}
\tightlist
\item
  \textbf{Log-Sum-Exp Trick:} To prevent catastrophic
  underflow/overflow:
  \[\log \sum_j \exp(x_j) = x_{\max} + \log \sum_j \exp(x_j - x_{\max})\]
\item
  \textbf{Preventing Label Switching:} We enforce strict ordering of the
  means (\(\mu_0 < \mu_1 < \dots\)) at each Gibbs iteration.
\item
  \textbf{Rcpp for Real-World Grids:} Recreated the local neighborhood
  and distance calculations in \textbf{C++ via Rcpp} to make iterations
  highly efficient.
\end{itemize}
\end{frame}

\section{4. Results \& Application}\label{results-application}

\subsection{Numerical Results}\label{numerical-results}

\begin{frame}{Simulation Results}
\phantomsection\label{simulation-results}
\begin{itemize}
\tightlist
\item
  \textbf{Parameter Recovery:} MWG algorithm shows rapid convergence
  after a short burn-in period.
\item
  \textbf{Credible Intervals:} True spatial and geographic parameters
  are captured within 90\% Bayesian credible intervals (5\% and 95\%
  quantiles).
\item
  \textbf{Geographic Trajectory:} Seed positions walk stochastically and
  converge exactly on their true coordinates.
\end{itemize}
\end{frame}

\begin{frame}{Image Segmentation Application}
\phantomsection\label{image-segmentation-application}
\begin{itemize}
\tightlist
\item
  \textbf{Texture Segmentation:} Without spatial priors, simple mixtures
  produce extremely noisy results.
\item
  \textbf{Our Model:} The latent field with seeds creates \textbf{smooth
  regions}, preserving concentric attraction near seeds.
\item
  \textbf{Visual Comparison:}

  \begin{enumerate}
  \tightlist
  \item
    Noisy Image \(Y\) \(\rightarrow\) heavy noise.
  \item
    Our Model \(\rightarrow\) clean, mathematically consistent shapes.
  \end{enumerate}
\end{itemize}
\end{frame}

\subsection{Conclusion}\label{conclusion}

\begin{frame}{Conclusion and Next Steps}
\phantomsection\label{conclusion-and-next-steps}
\begin{itemize}
\tightlist
\item
  \textbf{Key Conclusions:} The MWG pipeline makes inference in
  geographically complex HMRFs viable and robust.
\item
  \textbf{Efficiency:} Estimation via Pseudoposterior balances
  computational cost and statistical precision.
\item
  \textbf{Future Work:}

  \begin{itemize}
  \tightlist
  \item
    IDK
  \end{itemize}
\end{itemize}
\end{frame}




\end{document}
